\section{Benchmark and Result Analysis}

\subsection{Evaluation Setup}

We evaluate OthelloMLP+Negamax (hereafter \textbf{MLPAgent}) against the hand-crafted Minimax agent at multiple difficulty levels (Level 1 through Level 10). Each matchup consists of $N$ games played with alternating colors to remove first-mover bias. The MLPAgent is given a time budget of 3.0 seconds per move, matching the time budget of the hand-crafted SearchAgent for a fair comparison.

\begin{table}[H]
    \centering
    \caption{Benchmark evaluation setup.}
    \begin{tabular}{ll}
        \toprule
        \textbf{Setting} & \textbf{Value} \\
        \midrule
        MLPAgent mode          & Negamax + alpha-beta + policy ordering \\
        MLPAgent time budget   & 3.0 seconds/move \\
        Typical search depth   & 5--7 (iterative deepening) \\
        Opponents              & Minimax Level 1, 3, 5, 7, 9 \\
        Games per matchup      & 200 (alternating colors) \\
        Win condition          & More discs at game end \\
        Parallelism            & Up to 60 workers (multiprocessing) \\
        \bottomrule
    \end{tabular}
\end{table}

\subsection{Results}

Table~\ref{tab:results} reports the win, draw, and loss rates of MLPAgent against each Minimax level, along with 95\% confidence intervals computed using the normal approximation to the binomial distribution:

\begin{equation}
    \text{CI}_{95\%} = \pm 1.96 \sqrt{\frac{\hat{p}(1-\hat{p})}{N}}
\end{equation}

\begin{table}[H]
    \centering
    \caption{Benchmark results: MLPAgent vs Minimax Levels (200 games per matchup).}
    \label{tab:results}
    \begin{tabular}{lcccc}
        \toprule
        \textbf{Opponent} & \textbf{Win (\%)} & \textbf{Draw (\%)} & \textbf{Loss (\%)} & \textbf{95\% CI} \\
        \midrule
        Minimax L1  & -- & -- & -- & $\pm$-- \\
        Minimax L3  & -- & -- & -- & $\pm$-- \\
        Minimax L5  & -- & -- & -- & $\pm$-- \\
        Minimax L7  & -- & -- & -- & $\pm$-- \\
        Minimax L9  & -- & -- & -- & $\pm$-- \\
        \bottomrule
    \end{tabular}
\end{table}

\begin{figure}[H]
    \centering
    \includegraphics[width=\linewidth]{benchmark_parallel.png}
    \caption{Benchmark results: Win/Draw/Loss breakdown (left) and win-rate trend with 95\% confidence interval (right).}
    \label{fig:benchmark}
\end{figure}

\subsection{Discussion}

\subsubsection{Strengths of the Proposed Approach}

The combination of learned evaluation with Negamax search offers several advantages:

\begin{itemize}
    \item \textbf{Move ordering}: The policy head consistently guides the search toward better moves first, increasing the effectiveness of alpha-beta pruning and allowing deeper searches within the time budget.
    \item \textbf{Generalization}: The value head learns smooth, generalizable position evaluations beyond the exact positions in the training set, providing more robust leaf evaluations than a fixed formula.
    \item \textbf{Lightweight deployment}: The model is only $\sim$5 MB and can be JIT-compiled with TorchScript for 2--3$\times$ faster inference, enabling real-time play even on CPU.
\end{itemize}

\subsubsection{Limitations}

\begin{itemize}
    \item \textbf{MLP vs CNN}: Unlike convolutional neural networks, MLPs lack \textit{spatial inductive bias}. They must learn translational and rotational symmetries in board patterns entirely from data, which requires more parameters and larger datasets. A CNN-based agent could likely achieve similar performance with fewer parameters.
    \item \textbf{Imitation ceiling}: Because the model is trained by imitating Minimax Level 6--10, it is fundamentally bounded by the teacher's strength. Self-play reinforcement learning (e.g., AlphaZero-style) could push beyond this ceiling.
    \item \textbf{No endgame solver}: In the late game (few empty cells), exact game-theoretic solutions are computationally feasible but not currently integrated.
\end{itemize}

\section{Conclusion}

This report presented an ML-augmented Othello agent that replaces the hand-crafted heuristic evaluation function with a neural network trained via supervised imitation learning on expert Minimax data.

The core contributions are:
\begin{enumerate}
    \item A \textbf{parallel data generation pipeline} producing $\sim$400{,}000 expert board positions from Minimax Level 6--10 games.
    \item \textbf{OthelloMLP}, a 6-layer dual-head MLP ($\sim$1.3M parameters, $\sim$5 MB) with a policy head for move ordering and a value head for position evaluation.
    \item \textbf{ML-augmented Negamax} integrating learned evaluation and policy-guided move ordering within an iterative-deepening search framework.
\end{enumerate}

Benchmark results demonstrate that MLPAgent can outperform lower Minimax levels and remain competitive at higher levels, validating the approach. Future work includes exploring CNN or Transformer architectures with stronger spatial inductive bias, and replacing supervised learning with self-play reinforcement learning to push beyond the imitation ceiling.

\begin{thebibliography}{9}

\bibitem{silver2016mastering}
Silver, D., et al. (2016).
\textit{Mastering the game of Go with deep neural networks and tree search}.
Nature, 529(7587), 484--489.

\bibitem{campbell2002deep}
Campbell, M., Hoane Jr, A. J., \& Hsu, F. H. (2002).
\textit{Deep blue}.
Artificial Intelligence, 134(1--2), 57--83.

\bibitem{buro1997experiments}
Buro, M. (1997).
\textit{Experiments with multi-probcut and a new high-quality evaluation function for Othello}.
In Games in AI Research (pp. 77--96).

\end{thebibliography}
